%% kapa: exam=1 round=25 year=2014 period=2 form=A subject=appraisal_law subject_ko=부동산관계법규 seq=13 number=13
%% source: Q-Net 감정평가사 기출문제 (cst00309) artlSeq=5101207 file=extracted/1008/5101207/2014년 제25회 감정평가사 1차시험 기출문제/2교시(A형).hwp
%% conversion: kordoc 4.24.51 -> tools/gen1.py (auto); review=proofread-sonnet; flags=none
\begin{question}{13}
국토의 계획 및 이용에 관한 법령상 아래의 내용이 설명하는 것은?
\begin{framedtext}
시장ㆍ군수 또는 구청장은 토지거래계약에 관한 허가 신청이 있는 경우 공익사업용 토지에 대하여 국가, 지방자치단체, 한국토지주택공사, 그 밖에 대통령령으로 정하는 공공기관 또는 공공단체가 그 매수를 원하는 경우에는 이들 중에서 해당 토지를 매수할 자를 지정하여 그 토지를 협의 매수하게 할 수 있다. \par
\end{framedtext}
\begin{choices}
  \item 선매
  \item 수용
  \item 매수청구
  \item 환매
  \item 공용취득
\end{choices}
\end{question}
\answer{1}
